\documentclass{article}
\usepackage{amsmath}
\usepackage{tikz}
\usetikzlibrary{arrows.meta, decorations.markings,patterns}
\begin{document}


2.3-4

(1):the displacement field is uniform across both dielectrics:  

$$D = \sigma\qquad E_1 = \frac{D}{\varepsilon_1\varepsilon_0} = \frac{\sigma}{\varepsilon_1\varepsilon_0}, \quad E_2 = \frac{\sigma}{\varepsilon_2\varepsilon_0}.$$


$$U = E_1 d_1 + E_2 d_2 = \sigma \left( \frac{d_1}{\varepsilon_1\varepsilon_0} + \frac{d_2}{\varepsilon_2\varepsilon_0} \right).$$

 $Q = \sigma S$, so:  

$$C = \frac{Q}{U} = \frac{\sigma S}{\sigma\left( \frac{d_1}{\varepsilon_1\varepsilon_0} + \frac{d_2}{\varepsilon_2\varepsilon_0} \right)}= \frac{S}{\frac{d_1}{\varepsilon_1\varepsilon_0} + \frac{d_2}{\varepsilon_2\varepsilon_0}}.$$

$$C=\frac{\varepsilon_1\varepsilon_2\varepsilon_0S}{\varepsilon_1d_2+\varepsilon_2d_1}$$

(2): 
$$\sigma=\sigma_{e0}$$
$$
P_1 = D - \varepsilon_0 E_1 = \sigma_{e0} \left( 1 - \frac{1}{\varepsilon_1} \right),
$$
$$
P_2 = \sigma_{e0} \left( 1 - \frac{1}{\varepsilon_2} \right).
$$

$$\sigma_{p1} = P_1\quad  \sigma_{p2} = -P_2\qquad \sigma_e' = \sigma_{p1} + \sigma_{p2} = P_1 - P_2.$$



$$\sigma_e' = \sigma_{e0} \left[ \left(1 - \frac{1}{\varepsilon_1}\right) - \left(1 - \frac{1}{\varepsilon_2}\right) \right]= \sigma_{e0}  \left( \frac{1}{\varepsilon_2} - \frac{1}{\varepsilon_1} \right).$$

$$\sigma_e' = \frac{\varepsilon_1-\varepsilon_2}{\varepsilon_1\varepsilon_2}\sigma_{e0}$$



(3): From the result in part (1):  

$$U = \sigma_{e0} \left( \frac{d_1}{\varepsilon_1\varepsilon_0} + \frac{d_2}{\varepsilon_2\varepsilon_0} \right)=\frac{\left(\varepsilon_1d_2+\varepsilon_2d_1\right)\sigma_{e0}}{\varepsilon_1\varepsilon_2\varepsilon_0}$$

 (4): From Gauss’s law, $D$ is uniform across the dielectrics:  

$$D_1 = D_2 = \sigma_{e0}$$




\newpage


2.3-7

(1) From continuity of $ D $ and the voltage condition:

$$D = \frac{U}{\dfrac{d-t}{\varepsilon_0} + \dfrac{t}{\varepsilon\varepsilon_0}}.$$

Inside the dielectric:

$$
E_m = \frac{D}{\varepsilon\varepsilon_0} = \frac{U}{\varepsilon\varepsilon_0 \left( \dfrac{d-t}{\varepsilon_0} + \dfrac{t}{\varepsilon\varepsilon_0} \right)},\quad
P = D - \varepsilon_0 E_m = D\left( 1 - \frac{1}{\varepsilon} \right).
$$

$$E=\frac{U}{\varepsilon d+\left(1-\varepsilon\right)t},D=\frac{\varepsilon\varepsilon_0U}{\varepsilon d+\left(1-\varepsilon\right)t},P=\frac{\left(\varepsilon-1\right)\varepsilon_0U}{\varepsilon d+\left(1-\varepsilon\right)t}$$

(2) Charge $ Q $ on the plates

$$Q = D S = \frac{\varepsilon\varepsilon_0\boldsymbol{s}U}{\varepsilon d+(1-\varepsilon)t}$$

(3) Electric field $ E $ in the air gaps

$$E_a = \frac{D}{\varepsilon_0} = \frac{\varepsilon U}{\varepsilon d+(1-\varepsilon)t}$$

(4) Capacitance $ C $

$$C = \frac{Q}{U} =\frac{\varepsilon\varepsilon_0S}{\varepsilon d+(1-\varepsilon)t}$$


\newpage

2.3-12

   $$C = \frac{\varepsilon\varepsilon_0 S}{d}.\qquad C_1 = \dfrac{\varepsilon_1\varepsilon_0 S_1}{d} \quad C_2 = \dfrac{\varepsilon_2\varepsilon_0 S_2}{d} $$


   Since the two regions are in parallel (same voltage):

   $$C = C_1 + C_2 = \frac{\varepsilon_1\varepsilon_0 S_1}{d} + \frac{\varepsilon_2\varepsilon_0 S_2}{d}.$$

  $$ C = \frac{\varepsilon_0(\varepsilon_1 S_1 + \varepsilon_2 S_2)}{d} $$


\newpage

2.3-13


(1) 

dielectric area: area $ S/2 $, dielectric thickness $ t $, air thickness $ d-t $.  

  $$  C_A = \frac{S/2}{\dfrac{d-t}{\varepsilon_0} + \dfrac{t}{\varepsilon\varepsilon_0}}.$$

air area: area $ S/2 $, separation $ d $

  $$C_B = \frac{\varepsilon_0 S}{2d}.$$

Total capacitance (parallel):

$$C = C_A + C_B = \frac{S}{2} \left[ \frac{1}{\dfrac{d-t}{\varepsilon_0} + \dfrac{t}{\varepsilon\varepsilon_0}} + \frac{\varepsilon_0}{d} \right].$$

$$U = \frac{Q}{C} = \frac{Q}{\dfrac{S}{2} \left[ \dfrac{1}{\dfrac{d-t}{\varepsilon_0} + \dfrac{t}{\varepsilon\varepsilon_0}} + \dfrac{\varepsilon_0}{d} \right]}$$

$$U=\frac{2[\varepsilon d+(1-\varepsilon)t]Qd}{\varepsilon_0S[2\varepsilon d+(1-\varepsilon)t]}$$


(2) Capacitance $ C $

$$C = \frac{\varepsilon_0S[2\varepsilon d+(1-\varepsilon)t]}{2[\varepsilon d+(1-\varepsilon)t]d} $$



(3) Polarization charge surface density $ \sigma_e' $ on the dielectric

dielectric area:
$$D = \frac{Q_A}{S/2} = \frac{C_A U}{S/2}.$$

$$P = D \left( 1 - \frac{1}{\varepsilon} \right).$$

$$\sigma_e' = P = D \left( 1 - \frac{1}{\varepsilon} \right).\qquad D = \dfrac{Q}{S} \cdot \dfrac{2C_A}{C_A + C_B} $$

$$\sigma_e' = \frac{Q}{S} \cdot \frac{2C_A}{C_A + C_B} \left( 1 - \frac{1}{\varepsilon} \right)=\frac{2(\varepsilon-1)Qd}{S[2\varepsilon d+(1-\varepsilon)t]}$$






\newpage

2.3-19



(1)Gauss’s law:
$$D = \frac{Q}{4\pi r^2} \quad (r > R).$$

$$ R < r < a :E = \frac{D}{\varepsilon_0} = \frac{Q}{4\pi\varepsilon_0 r^2}.$$

$$ a < r < b :E = \frac{D}{\varepsilon\varepsilon_0} = \frac{Q}{4\pi\varepsilon\varepsilon_0 r^2}.$$

$$ r > b :E = \frac{Q}{4\pi\varepsilon_0 r^2}.$$

\vspace{1.5em}

(2) Polarization $ P $ and surface polarization charge density $ \sigma_e' $

$$P = D - \varepsilon_0 E = \frac{Q}{4\pi r^2} \left( 1 - \frac{1}{\varepsilon} \right).$$

Surface polarization charge density ($ \sigma_p' = \mathbf{P} \cdot \mathbf{n}_{\text{out}} $):

$$\sigma_e'(a) = -P(a) = -\frac{Q}{4\pi a^2} \left( 1 - \frac{1}{\varepsilon} \right).$$

$$\sigma_e'(b) = P(b) = \frac{Q}{4\pi b^2} \left( 1 - \frac{1}{\varepsilon} \right).$$

$$P=\frac{\left(\varepsilon-1\right)Q}{4\pi\varepsilon r^{2}},\sigma_{e}'(a)=-\frac{\left(\varepsilon-1\right)Q}{4\pi\varepsilon a^{2}},\sigma_{e}'(b)=\frac{\left(\varepsilon-1\right)Q}{4\pi\varepsilon b^{2}}$$

\vspace{2em}

(3) Volume polarization charge density $ \rho_e' $ inside the dielectric

$$\rho_e' = -\nabla \cdot \mathbf{P}.$$

Since $ P(r) \propto 1/r^2 $, we have $ \nabla \cdot \mathbf{P} = 0 $, hence:

$$\rho_e' = 0$$






\newpage

2.3-24

By Gauss’s law:

$$D(\rho) \cdot 2\pi\rho l = Q \quad \Rightarrow \quad D(\rho) = \frac{Q}{2\pi\rho l}.$$

Electric field:

$$E_1(\rho) = \frac{D}{\varepsilon_1\varepsilon_0} = \frac{Q}{2\pi\varepsilon_1\varepsilon_0 l \rho} \quad (a < \rho < r),$$

$$E_2(\rho) = \frac{Q}{2\pi\varepsilon_2\varepsilon_0 l \rho} \quad (r < \rho < b).$$

Voltage:

$$U = \int_a^r E_1 \, d\rho + \int_r^b E_2 \, d\rho= \frac{Q}{2\pi l} \left[ \frac{1}{\varepsilon_1\varepsilon_0} \ln\frac{r}{a} + \frac{1}{\varepsilon_2\varepsilon_0} \ln\frac{b}{r} \right].$$

Capacitance:

$$C = \frac{Q}{U} = \frac{2\pi\varepsilon_1\varepsilon_2\varepsilon_0l}{\varepsilon_2\ln(\frac{r}{a})+\varepsilon_1\ln(\frac{b}{r})}$$


\newpage

2.3-35


By Gauss’s law:$D(\rho) = \frac{\lambda}{2\pi\rho}$

$$E_in(\rho) = \frac{D}{\varepsilon_1\varepsilon_0} = \frac{\lambda}{2\pi\varepsilon_1\varepsilon_0 \rho} \quad (R_1 < \rho < r)$$ 
$$E_out(\rho) = \frac{D}{\varepsilon_2\varepsilon_0} = \frac{\lambda}{\pi\varepsilon_1\varepsilon_0 \rho} \quad (r < \rho < R_2)$$

$$\rho = R_1^+ ,\quad E_{in,\max} = \frac{\lambda}{2\pi\varepsilon_1\varepsilon_0 R_1}$$
$$\rho = r^+ ,\quad E_{out,\max} = \frac{\lambda}{\pi\varepsilon_1\varepsilon_0 r}$$

$$ E_{out,\max} = 2 \cdot \frac{R_1}{r} \cdot E_{in,\max} > E_{in,\max} $$

$$ \rm{while}\quad  E_{out,\max}=E_m, E_{in,\max}<E_m $$

\vspace{2em}

\textbf{So the outer dielectric breaks down first.}



$$\lambda = \pi\varepsilon_1\varepsilon_0 r E_m ,E_1(\rho) = \frac{r E_m}{2\rho}, E_2(\rho) = \frac{r E_m}{\rho}$$


$$U_m = \int_{R_1}^{r} \frac{r E_m}{2\rho} \, d\rho + \int_{r}^{R_2} \frac{r E_m}{\rho} \, d\rho= \frac{r E_m}{2} \ln\frac{r}{R_1} + r E_m \ln\frac{R_2}{r}.$$

\vspace{2em}

$$U_m = \frac{r E_m}{2} \left[ \ln\frac{r}{R_1} + 2\ln\frac{R_2}{r} \right]= \frac{r E_m}{2} \ln\frac{R_2^2}{r R_1}.$$






\newpage

2.3-39

Before 

$$C_0 = \frac{\varepsilon_0 S}{d}, \quad Q_0 = C_0 U, \quad W_0 = \frac{1}{2} C_0 U^2$$

After 

$$C = \frac{\varepsilon\varepsilon_0 S}{d}, \quad Q = C U, \quad W = \frac{1}{2} C U^2$$



(1) Change in electrostatic energy
$$\Delta W_e = W - W_0 = \frac{(\varepsilon\varepsilon_0 - \varepsilon_0)S U^2}{2d}=\frac{(\varepsilon-1)\varepsilon_0SU^2}{2d}$$




(2) 
$$\Delta Q = Q - Q_0 = \frac{(\varepsilon\varepsilon_0 - \varepsilon_0)S U}{d} $$

$$W_{\text{source}} = U \cdot \Delta Q = \frac{(\varepsilon\varepsilon_0 - \varepsilon_0)S U^2}{d}=\frac{(\varepsilon-1)\varepsilon_0SU^2}{d} $$




(3) 
$$W_{\text{source}} = \Delta W_e + W_{\text{diel}}$$

$$W_{\text{diel}} = W_{\text{source}} - \Delta W_e= \frac{(\varepsilon\varepsilon_0 - \varepsilon_0)S U^2}{2d}=\frac{(\varepsilon-1)\varepsilon_0SU^2}{2d} $$





\end{document}